\section*{Planteamiento del Problema}
\addcontentsline{toc}{section}{Planteamiento del Problema}

A diferencia del diseño digital, que se apoya en abstracciones estandarizadas y flujos de automatización maduros, el diseño de circuitos analógicos descansa en el conocimiento experto, el razonamiento heurístico y un ajuste iterativo guiado por simulación \cite{analogagent}. Ese ciclo ---proponer un circuito, simularlo, revisar los resultados y volver a ajustar--- se repite numerosas veces antes de obtener un diseño que cumpla las metas, y validar un diseño completo puede exigir largos periodos de iteración \cite{pcbschemagen}. Buena parte de ese conocimiento reside en la experiencia de quien diseña y no en un procedimiento que pueda seguirse paso a paso, lo que vuelve el diseño analógico difícil tanto de automatizar como de aprender.

Los Modelos de Lenguaje Grande (LLM) abrieron la posibilidad de describir un circuito en lenguaje natural y recibir una propuesta de diseño de forma automática. Al llevarla a la práctica, sin embargo, aparecen limitaciones documentadas. Estos modelos alucinan componentes, violan restricciones físicas del circuito y producen salidas eléctricamente inconsistentes \cite{circuitlm}, por lo que sus diseños rara vez funcionan en el primer intento, ya que el modelo no comprueba por sí mismo si lo que generó cumple. La respuesta de la investigación ante esto ha sido sustituir el modelo único por varios agentes que se reparten el trabajo, acompañados de un mecanismo que verifica cada propuesta antes de darla por válida.

Los trabajos publicados se reparten entre dos líneas. La primera diseña a nivel de transistor, atado a un proceso de fabricación específico descrito por su \textit{Process Design Kit} (PDK), y emplea la simulación SPICE como realimentación durante el diseño; es la línea más desarrollada \cite{analogcoder, analogagent, analogsage}. La segunda trabaja a nivel de placa, con componentes comerciales, y verifica el resultado mediante reglas eléctricas \cite{circuitlm, pcbschemagen}, en parte porque no todo componente comercial cuenta con un modelo SPICE utilizable, lo que impide simular cualquier diseño de forma uniforme. La simulación SPICE a nivel de placa ha comenzado a aparecer, pero en forma de banco de prueba que compara el diseño generado contra un circuito de referencia construido por expertos \cite{hwebench}, y no como una señal que el propio sistema utilice para ajustar el diseño durante su ejecución. Entre los trabajos revisados no se encontró un sistema generativo de varios agentes que use la simulación SPICE como realimentación dentro del lazo de diseño, con un curador que decida el ajuste, orientado a circuitos analógicos CA/CD con componentes comerciales. Ese es el punto del espacio de diseño que este proyecto se propone ocupar.

De lo anterior se desprende la pregunta que ordena el trabajo:

\begin{quote}
    \textit{¿Es posible construir un ecosistema de varios agentes que, a partir de una descripción en lenguaje natural, genere un circuito analógico CA/CD con componentes comerciales, lo simule en ngspice y ajuste sus valores de forma automática hasta cumplir las metas de diseño? ¿Y cómo se comporta su desempeño frente al de un enfoque de un solo modelo que repite el intento varias veces?}
\end{quote}
